\section{Conclusion}
\label{sec:conclusion}

We began with a counterintuitive bet:
that a neural network encoder trained on nothing but small MLP and CNN checkpoints
could predict ViT model accuracy better than the state of the art ---
without ever seeing a transformer.
Our experiments confirm the bet pays off, and the mechanism is both simpler and more
surprising than we anticipated.
It is the coordinate system, not the symmetry group.

Graph-based SSL with permutation equivariance achieves $R^2 = 0.231$ on 53 held-out
ViT-S/16 ImageNet models, versus $R^2 = 0.072$ for SANE's flat tokenization ---
a ${\sim}220\%$ improvement with no ViT training data.
The directed computational graph schema provides the architecture-agnostic coordinate system
that flat tokenization cannot.
The secondary contribution is a principled negative result:
scale+permutation equivariance does not improve over permutation-only in the ViT
cross-architecture SSL setting ($\Delta R^2 = -0.021$ in the controlled h-m1 evaluation;
$\Delta R^2 = -0.143$ in the h-m2 ablation, though partially confounded by training duration),
consistent with our LayerNorm gauge-fixing hypothesis.
The third contribution is methodological:
we identify the MMD confounding problem and recommend $R^2$-based evaluation for
cross-architecture property prediction.

\textbf{Future Directions.}
(1) \emph{Normalization-conditional equivariance:}
Test scale equivariance benefit on non-LayerNorm architectures to confirm the gauge-fixing
hypothesis.
(2) \emph{Multi-seed, full-zoo validation:}
5 seeds and 250+ ViT models for statistical significance and calibrated magnitude estimates.
(3) \emph{Functional decoder design:}
Hypernetwork-style decoder with architecture-aware output projections to enable latent
interpolation.
(4) \emph{Controlled ablation:}
Re-run EquiSSL training for 100 epochs to provide an epoch-matched ablation comparison.

In a field that increasingly treats model checkpoints as data,
the ability to analyze an arbitrary new architecture using only existing model zoo data is
a fundamental primitive.
This work provides initial proof-of-concept evidence that the primitive exists ---
the coordinate system is the key ---
and opens a research agenda for understanding exactly when and how it works.
